% Part: proof-theory
% Chapter: natural-deduction
% Section: rules-N1

\documentclass[../../../include/open-logic-section]{subfiles}

\begin{document}

\begin{table}
\begin{center}
\begin{tabular}[t]{@{}cc@{}}
\bottomAlignProof
\AxiomC{$\Discharge{!A}{x}$}
\shortDeduce
\DeduceC{$!B$}
\DischargeRule{\Intro{\lif}}{x}
\UnaryInfC{$!A \lif !B$}
\DisplayProof
&
\bottomAlignProof
\AxiomC{$!A \lif !B$}
\AxiomC{$!A$}
\RightLabel{\Elim{\lif}}
\BinaryInfC{$!B$}
\DisplayProof
\\[4ex]
\bottomAlignProof
\AxiomC{$!A$}
\AxiomC{$!B$}
\RightLabel{\Intro{\land}}
\BinaryInfC{$!A \land !B$}
\DisplayProof
&
\begin{tabular}[b]{@{}l@{}}
\AxiomC{$!A \land !B$}
\RightLabel{\Elim{\land}[1]}
\UnaryInfC{$!A$}
\DisplayProof
\\[3ex]
\bottomAlignProof
\AxiomC{$!A \land !B$}
\RightLabel{\Elim{\land}[2]}
\UnaryInfC{$!B$}
\DisplayProof
\end{tabular}
\\[4ex]
\begin{tabular}{@{}l@{}}
\AxiomC{$!A$}
\RightLabel{\Intro{\lor}[1]}
\UnaryInfC{$!A \lor !B$}
\DisplayProof
\\[3ex]
\AxiomC{$!B$}
\RightLabel{\Intro{\lor}[2]}
\UnaryInfC{$!A \lor !B$}
\DisplayProof
\end{tabular}
&
\AxiomC{$!A \lor !B$}
\AxiomC{$\Discharge{!A}{x}$}
\shortDeduce
\DeduceC{$!C$}
\AxiomC{$\Discharge{!B}{y}$}
\shortDeduce
\DeduceC{$!C$}
\DischargeRule{\Elim{\lor}}{x\,y}
\TrinaryInfC{$!C$}
\DisplayProof
\bottomAlignProof
\\[6ex]
\bottomAlignProof
\AxiomC{$\Discharge{!A}{x}$}
\shortDeduce
\DeduceC{$\lfalse$}
\DischargeRule{\Intro{\lnot}}{x}
\UnaryInfC{$\lnot !A$}
\DisplayProof
&
\bottomAlignProof
\AxiomC{$\lnot !A$}
\AxiomC{$!A$}
\RightLabel{\Elim{\lnot}}
\BinaryInfC{$\lfalse$}
\DisplayProof
\\[4ex]
\bottomAlignProof
\AxiomC{$!A(a)$}
\RightLabel{\Intro{\lforall}}
\UnaryInfC{$\lforall[x][!A(x)]$}
\DisplayProof
&
\bottomAlignProof
\AxiomC{$\lforall[x][!A(x)]$}
\RightLabel{\Elim{\lforall}}
\UnaryInfC{$!A(t)$}
\DisplayProof
\\[4ex]
\bottomAlignProof
\AxiomC{$!A(t)$}
\RightLabel{\Intro{\lexists}}
\UnaryInfC{$\lexists[x][!A(x)]$}
\DisplayProof
&
\bottomAlignProof
\AxiomC{$\lexists[x][!A(x)]$}
\AxiomC{$\Discharge{!A(a)}{x}$}
\shortDeduce
\DeduceC{$!C$}
\DischargeRule{\Elim{\lexists}}{x}
\BinaryInfC{$!C$}
\DisplayProof
\\[4ex]
\bottomAlignProof
\AxiomC{$\lfalse$}
\RightLabel{\FalseInt}
\UnaryInfC{$!A$}
\DisplayProof
&
\bottomAlignProof
\AxiomC{$\Discharge{\lnot!A}{x}$}
\shortDeduce
\DeduceC{$\lfalse$}
\DischargeRule{\FalseCl}{x}
\UnaryInfC{$!A$}
\DisplayProof
\end{tabular}
\end{center}
\caption{\Log{N1i} ও \Log{N1c}-র বিধি। $\Intro{\lforall}$-এ
!!{constant}~$a$ সিদ্ধান্ত বা কোনো খোলা অনুমানে ঘটে না।
$\Elim{\lexists}$-এ $a$ প্রধান অস্তিত্বসূচক পূর্বধারণা,
শেষ সিদ্ধান্ত, কিংবা ওই বিধির বিশেষ অব্যাহতি-পাওয়া
অনুমান ছাড়া অন্য কোনো খোলা অনুমানে ঘটে না।
\Log{N1i}-তে \FalseCl নেই।}
% BN-SRC-547 TARGET-CORRECTION-BEGIN
\par\smallskip\noindent\textnormal{\small\textbf{উৎস-সংশোধন (BN-SRC-547):} অস্তিত্ব-নিষ্কাশনে আইগেনচলটি বিশেষ অব্যাহতি-পাওয়া গৌণ অনুমানে থাকতেই পারে; উৎসের ``কোনো পূর্বধারণায় নয়'' বক্তব্যের বদলে প্রধান অস্তিত্বসূচক পূর্বধারণা ও অন্য খোলা অনুমানগুলি থেকে অনুপস্থিতি বলা হয়েছে।}
% BN-SRC-547 TARGET-CORRECTION-END
\ollabel{tab:N1}
\end{table}

\end{document}
